% ==================================================================
% id:       app-004
% topic:    שימושי הנגזרת
% subtopic: חקירת פונקציה
% solution: full
% ==================================================================
\textbf{מצאו את תחומי העליה והירידה ואת נקודות הקיצון המקומי של הפונקציות הבאות:}
\bE\bMs{2}
\item $\ds{ f(x) = \frac{2x}{1 + x^2}} $
\begin{solbox}
נגזור
$$.f'(x)=\frac{2(1-x^2)}{(1+x^2)^2}=\frac{2(1-x)(1+x)}{(1+x^2)^2}$$
מתקיים
$f'(x)=0$
אמ"מ
$x=1$
או
$x=-1$.
נשים לב כי
$f'(x)>0$
אמ"מ
$|x|<1$
ו-
$f'(x)<0$
אמ"מ
$|x|>1$,
לכן בתחום
$|x|<1$
הפונקציה עולה ובתחום
$|x|>1$
הפונקציה יורדת.
מכאן נובע בפרט כי
$x=1$
היא נקודת מקסימום מקומי ואילו
$x=-1$
היא נקודת מינימום מקומי.
\end{solbox}
\item $\ds{f(x) = x \cdot \sqrt[3]{x - 1}}$
\hint{שימו לב לנקודה שבה הפונקציה אינה גזירה: גם בה צריך לבדוק את סימן הנגזרת משני הצדדים.}
\begin{solbox}
נגזור ונשווה ל-
$0$:
$$f'(x)=\frac{4x-3}{3(x-1)^{\frac{2}{3}}}=0$$
אמ"מ
$x=\frac{3}{4}$.
נשים לב כי
$f'(x)<0$
אמ"מ
$x<\frac{3}{4}$
ו-
$f'(x)>0$
אמ"מ
$x>\frac{3}{4}$
כאשר
$x\neq 1$
(הפונקציה לא גזירה בנקודה זו),
לכן בתחום
$x<\frac{3}{4}$
הפונקציה יורדת ובתחום
$x>\frac{3}{4}$
הפונקציה עולה.
מכאן נובע בפרט כי
$x=\frac{3}{4}$
היא נקודת מינימום מקומי.
כאמור,
$x=1$
הפונקציה לא גזירה. יתכן מאוד ויש לה קיצון בנקודה זו, אך מחקירת תחומי העליה וירידה של הפונקציה הנתונה עולה כי
$x=1$
אינה נקודת קיצון של
$f$.
\end{solbox}
\item $\ds{f(x) = x e^{-x}}$
\begin{solbox}
נגזור ונשווה ל-
$0$:
$$f'(x)=e^{-x}(1-x)=0$$
אמ"מ
$x=1$.
נשים לב כי
$f'(x)>0$
אמ"מ
$x<1$
ו-
$f'(x)<0$
אמ"מ
$x>1$,
לכן בתחום
$x<1$
הפונקציה עולה ובתחום
$x>1$
הפונקציה יורדת.
מכאן נובע בפרט כי
$x=1$
היא נקודת מקסימום מקומי.
\end{solbox}

\item $\ds{f(x) = \frac{\ln^2{x}}{x}}$
\begin{solbox}
נגזור ונשווה ל-
$0$:
$$f'(x)=\frac{\ln x\,(2-\ln x)}{x^2}=0$$
אמ"מ
$x=1$
או
$x=e^2$
נשים לב כי
$f'(x)<0$
אמ"מ
$0<x<1$
או
$x>e^2$
ו-
$f'(x)>0$
אמ"מ
$1<x<e^2$,
לכן בתחום
$(0,1)\cup(e^2,\infty)$
הפונקציה יורדת ובתחום
$(1,e^2)$
הפונקציה עולה.
מכאן נובע בפרט כי
$x=1$
היא נקודת מינימום מקומי
ו-
$x=e^2$
היא נקודת מקסימום מקומי.
\end{solbox}
\item $\ds{f(x) = \arctan{x} - \frac{1}{2} \ln(1 + x^2)}$
\begin{solbox}
נגזור ונשווה ל-
$0$:
$$f'(x)=\frac{1-x}{1+x^2}=0$$
אמ"מ
$x=1$.
נשים לב כי
$f'(x)>0$
אמ"מ
$x<1$
ו-
$f'(x)<0$
אמ"מ
$x>1$,
לכן בתחום
$x<1$
הפונקציה עולה ובתחום
$x>1$
הפונקציה יורדת.
מכאן נובע בפרט כי
$x=1$
היא נקודת מקסימום מקומי.
\end{solbox}

\item $\ds{f(x) = \cos{x} + \frac{1}{2} \cos{2x}}$
\hint{הפונקציה מחזורית. פרקו את הנגזרת לגורמים בעזרת $\sin2x=2\sin x\cos x$ וחקרו את סימנה בתוך מחזור אחד.}
\begin{solbox}
נגזור ונשווה ל-
$0$:
$$f'(x)=-\sin x-\sin(2x)=-\sin x(1+2\cos x)=0$$
אמ"מ
$x=\pi k$
או
$x=\frac{2\pi}{3}+2\pi l$
או
$x=\frac{4\pi}{3}+2\pi m$
כאשר
$k,l,m\in\mathbb{Z}$.
\iffalse
הפעם, לשם שינוי, נשתמש במבחן הנגזרת השנייה.
מתקיים
$$.f''(x)=-\cos x-2\cos(2x)$$
מתקיים
$$,f''(\pi k)<0,\ f''(\tfrac{2\pi}{3}+2\pi l)>0,\ f''(\tfrac{4\pi}{3}+2\pi l)>0$$
לכן
$x=\pi k$
הן נקודות מקסימום מקומי
לכל
$k\in\mathbb{R}$
ו-
$x=\frac{2\pi}{3}+2\pi l$, $\frac{4\pi}{3}+2\pi m$
הן נקודות מינימום מקומי
לכל
$l,m\in\mathbb{R}$.
\fi
לקביעת תחומי עליה וירידה נצטמצם לקטע בטווח המחזור של הפונקציה, לדוגמא
$[0,2\pi]$
נשים לב כי
$$f'(x)=-\sin x(1+2\cos x)<0$$
אמ"מ
$0<x<\frac{2\pi}{3}$
או
$\pi<x<\frac{4\pi}{3}$.
כמוכן,
$$f'(x)=-\sin x(1+2\cos x)>0$$
אמ"מ
$\frac{2\pi}{3}<x<\pi$
או
$\frac{4\pi}{3}<x<2\pi$.
מכאן
$f$
יורדת ב-
$$\left(2\pi k,\frac{2\pi}{3}+2\pi k\right)\bigcup\left(\pi+2\pi k,\frac{4\pi}{3}+2\pi k\right),\ k\in\mathbb{Z}$$
ועולה ב-
$$.\left(\frac{2\pi}{3}+2\pi k,2\pi (k+1)\right)\bigcup\left(\frac{4\pi}{3}+2\pi k,\pi+2\pi (k+1)\right),\ k\in\mathbb{Z}$$
\end{solbox}
\eMs
\eE

%\item הוכיחו כי \( |3x - x^3| \leq 2 \) עבור \( |x| \leq 2 \).
